% main_prob_lecture5_metropolis.tex
\documentclass[aspectratio=169]{beamer}

% --- Language & encoding ---
\usepackage[utf8]{inputenc}
\usepackage[T1]{fontenc}
\usepackage[english,french]{babel}
\usepackage{lmodern}
\usepackage{microtype}

% --- Math & tables ---
\usepackage{amsmath,amssymb}
\usepackage{booktabs}
\usepackage{graphicx}
\usepackage{pgfplots}
\pgfplotsset{compat=1.18}
\usepackage{changepage}

% --- Metropolis Theme ---
\usetheme[numbering=fraction,progressbar=frametitle]{metropolis}
\setbeamertemplate{frame numbering}[fraction]

% --- Colors (same palette as prior lectures) ---
\definecolor{BootcampBlueDark}{HTML}{1A3C68}
\definecolor{TextBlack}{HTML}{000000}
\definecolor{AccentGold}{HTML}{DAA520}
\definecolor{Crimson}{HTML}{990000}
\definecolor{MatrixGreen}{HTML}{228B22}
\definecolor{MatrixRed}{HTML}{B22222}

\setbeamercolor{block title example}{bg=BootcampBlueDark!90!black,fg=white}
\setbeamercolor{block body example}{bg=BootcampBlueDark!5!white,fg=black}

% === Thick frame title band ===
\setbeamercolor{frametitle}{bg=BootcampBlueDark,fg=white}
\setbeamerfont{frametitle}{series=\bfseries,size=\large}
\setbeamertemplate{frametitle}{
  \nointerlineskip
  \begin{beamercolorbox}[wd=\paperwidth,ht=3.2ex,dp=1.4ex,leftskip=1em,rightskip=1em]{frametitle}
    \usebeamerfont{frametitle}\insertframetitle
  \end{beamercolorbox}
}

% === Custom Definition Block ===
\newenvironment{definitionblock}[1]{%
  \par\vspace{0.4em}%
  \begin{beamercolorbox}[
      wd=\textwidth,
      sep=0.4em,
      leftskip=0.3em,
      rightskip=0.6em,
      rounded=false,
      shadow=false
    ]{block title example}%
    \raisebox{0.15em}{\usebeamerfont{block title example}#1}%
  \end{beamercolorbox}%
  \vspace{-0.4em}%
  \begin{beamercolorbox}[
      wd=\textwidth,
      sep=1em,
      leftskip=0.6em,
      rightskip=0.6em,
      rounded=false,
      shadow=false
    ]{block body example}%
    \usebeamerfont{block body example}%
}{%
  \end{beamercolorbox}%
  \par\vspace{0.6em}%
}

% === Global text formatting ===
\setbeamercolor{normal text}{fg=TextBlack,bg=white}
\setbeamercolor{title}{fg=TextBlack}
\setbeamercolor{structure}{fg=BootcampBlueDark}
\setbeamercolor{progress bar}{fg=BootcampBlueDark,bg=gray!30}

% --- Course Info ---
\title{Random Vectors and Convergence of Random Variables}
\institute{NYU Tandon School, Finance and Risk Engineering}
\author{Nathan Sauldubois}
\date{\today}
% \titlegraphic{\includegraphics[height=1.2cm]{your_logo.pdf}}

\metroset{
  sectionpage=progressbar,
  subsectionpage=progressbar
}

% --- Section & subsection transitions (plain) ---
\AtBeginSection{
  \frame[plain]{\sectionpage}
}
\AtBeginSubsection{
  \frame[plain]{\subsectionpage}
}

\begin{document}

\maketitle

% ===== Outline =====
\begin{frame}{Outline}
  \tableofcontents[hideallsubsections]
\end{frame}

% ===========================
% SECTION 0 — Motivation
% ===========================
% ===========================
% SECTION — Motivation & Overview
% ===========================
\section{Motivation \& Overview}

\begin{frame}{Outline}
  \tableofcontents[currentsection, hideothersubsections]
\end{frame}

\subsection{Where We're Going}

\begin{frame}{From Random Variables to Random Vectors and Convergence}

\begin{itemize}
  \item Extend random variables to \textbf{random vectors} and joint distributions.
  \item Understand dependence through covariance, correlation, and joint laws.
  \item Introduce \textbf{modes of convergence} for random variables:
  \begin{itemize}
    \item almost sure convergence,
    \item convergence in probability,
    \item convergence in \(L^p\),
    \item convergence in distribution.
  \end{itemize}
  \item Clarify the relationships between these notions and their practical meaning.
  \item Prepare the ground for limit theorems and statistical applications.
\end{itemize}

\end{frame}
% ======================================
% SECTION 1 — Random Variables (cont.)
% ======================================
\section{Random Vectors}

\begin{frame}{Outline}
  \tableofcontents[currentsection, hideothersubsections]
\end{frame}

% ======================================
% Random Vectors
% ======================================

\subsection{Definition}

\begin{frame}{Random Vectors}
\begin{definitionblock}{Random Vector}
A \emph{random vector} is a vector of random variables
\[
X = (X_1,\dots,X_d) \in \mathbb{R}^d,
\]
defined on the same probability space \((\Omega,\mathbb{P})\).
\pause
\end{definitionblock}

\begin{itemize}
  \item Each component \(X_k\) is a real-valued random variable.\pause
  \item The distribution of \(X\) is a probability measure on \(\mathbb{R}^d\).\pause
  \item The distribution of each \(X_k\) is called a \emph{marginal distribution}.\pause
\end{itemize}

\end{frame}

\begin{frame}{Example of a Random Vector}

\textbf{Example — Two-asset returns}

Let \(R_1\) and \(R_2\) denote the (random) returns of two assets over one period.

Define the random vector:
$
R = (R_1, R_2) \in \mathbb{R}^2.
$

\begin{itemize}
  \item \(R_1\): return of asset 1,
  \item \(R_2\): return of asset 2.
\end{itemize}
\pause
\medskip

Each component has its own marginal distribution:
\[
R_1 \sim \mathcal{N}(\mu_1,\sigma_1^2),
\qquad
R_2 \sim \mathcal{N}(\mu_2,\sigma_2^2).
\]

\medskip

\textbf{Joint behavior:}
\begin{itemize}
  \item dependence is captured by \(\mathrm{Cov}(R_1,R_2)\),
  \item diversification effects depend on the joint distribution of \(R\).
\end{itemize}

\end{frame}


% ======================================
% Joint and Marginal Densities
% ======================================

\subsection{Joint \& Marginal Densities}

\begin{frame}{Joint and Marginal Densities}
\begin{definitionblock}{Joint Density}
A random vector \((X,Y)\) has a \emph{joint density} \(f_{X,Y}\) if
\[
\mathbb{P}((X,Y)\in A) = \int_A f_{X,Y}(x,y)\,dx\,dy.
\]\pause
\end{definitionblock}

\begin{itemize}
  \item Marginal densities are obtained by integration:
  \[
    f_X(x)=\int_{\mathbb{R}} f_{X,Y}(x,y)\,dy,
    \quad
    f_Y(y)=\int_{\mathbb{R}} f_{X,Y}(x,y)\,dx.
  \]
\end{itemize}
\end{frame}

\begin{frame}{Example: Joint and Marginal Densities}

\textbf{Example — Uniform distribution on a square}

Let \((X,Y)\) have joint density
\[
f_{X,Y}(x,y)=
\begin{cases}
1, & (x,y)\in[0,1]\times[0,1],\\
0, & \text{otherwise.}
\end{cases}
\vspace{-5mm}
\]\pause

\medskip

\textbf{Step 1: Compute the marginal density of \(X\).}
\vspace{-1mm}
\[
f_X(x)=\int_0^1 1\,dy
=
\begin{cases}
1, & x\in[0,1],\\
0, & \text{otherwise.}
\end{cases}
\vspace{-2mm}
\]
\pause
\textbf{Step 2: Compute the marginal density of \(Y\).}
\[
f_Y(y)=\int_0^1 1\,dx
=
\begin{cases}
1, & y\in[0,1],\\
0, & \text{otherwise.}
\end{cases}
\vspace{-9mm}
\]\pause

\medskip

\textbf{Conclusion:}
\[
f_{X,Y}(x,y)=f_X(x)f_Y(y),
\]
so \(X\) and \(Y\) are independent.
\end{frame}

\begin{frame}{Independence via Densities}
\begin{definitionblock}{Independence}
Random variables \(X\) and \(Y\) are independent if
\[
f_{X,Y}(x,y) = f_X(x)\,f_Y(y)
\quad \text{(a.e.)}.
\]\pause
\end{definitionblock}

\begin{itemize}
  \item Independence \(\Rightarrow\) zero covariance (but not conversely).\pause
  \item Independence implies factorization of:
  \begin{itemize}
    \item joint CDF,\pause
    \item joint density,\pause
    \item characteristic function.\pause
  \end{itemize}
\end{itemize}
\end{frame}

\begin{frame}{How to Characterize the Law of a Random Vector}

Let \(X=(X_1,\dots,X_d)\) be a random vector in \(\mathbb{R}^d\).
\pause
\medskip

The distribution of \(X\) can be characterized in several equivalent ways:
\begin{itemize}
  \item by its \textbf{joint cumulative distribution function (CDF)};\pause
  \item by its \textbf{joint density} (if it exists);\pause
  \item by its \textbf{moments} (mean vector, covariance matrix);\pause
  \item by its \textbf{characteristic function}.\pause
\end{itemize}

\medskip

Each representation emphasizes different aspects of the dependence structure.

\end{frame}

\begin{frame}{Characteristic Function of a Random Vector}

The characteristic function of \(X\in\mathbb{R}^d\) is
\[
\Phi_X(t)
=
\mathbb{E}\big[e^{i\,t^\top X}\big],
\quad t\in\mathbb{R}^d.
\]
\pause
\begin{itemize}
  \item \(\Phi_X\) always exists.\pause
  \item \(\Phi_X\) uniquely determines the distribution of \(X\).\pause
  \item Independence is equivalent to factorization:
  \[
    \Phi_X(t)=\prod_{i=1}^d \Phi_{X_i}(t_i).
  \]\pause
\end{itemize}

\medskip

\textbf{Widely used for:}
limit theorems, Gaussian vectors, convergence in distribution.

\end{frame}

\begin{frame}{Summary: Characterizing a Random Vector}

\begin{center}
\begin{tabular}{ll}
\toprule
Tool & What it captures \\
\midrule
Joint CDF & Full distribution (always) \\
Joint density & Local behavior (if exists) \\
Moments & Location and dependence \\
Characteristic function & Full distribution (always) \\
\bottomrule
\end{tabular}
\end{center}

\medskip

\textbf{Key message:}  
different tools, same distribution — choose the right one for the problem.

\end{frame}

\begin{frame}{Example: Characterization via Joint Density}

Let \((X,Y)\) be a random vector with joint density
\[
f_{X,Y}(x,y)
=
\begin{cases}
e^{-(x+y)}, & x\ge0,\ y\ge0,\\
0, & \text{otherwise.}
\end{cases}
\vspace{-5mm}
\]
\pause
\medskip

\textbf{Step 1: Marginal densities}
\[
f_X(x)=\int_0^\infty e^{-(x+y)}\,dy = e^{-x}\mathbf{1}_{x\ge0},
\quad
f_Y(y)=e^{-y}\mathbf{1}_{y\ge0}.
\vspace{-5mm}
\]
\pause
\medskip

\textbf{Step 2: Factorization}
\[
f_{X,Y}(x,y)=f_X(x)f_Y(y).
\vspace{-9mm}
\]
\pause
\medskip

\textbf{Conclusion:}
\begin{itemize}
  \item \(X\sim \mathrm{Exp}(1)\), \(Y\sim \mathrm{Exp}(1)\)
  \item \(X\) and \(Y\) are independent
  \item the joint law is fully characterized by the density
\end{itemize}

\end{frame}



% ======================================
% Covariance and Correlation
% ======================================

\subsection{Covariance and Correlation}


\begin{frame}{Mean of a Random Vector}

\begin{definitionblock}{Mean Vector}
Let \(X=(X_1,\dots,X_d)\) be a random vector such that
\(\mathbb{E}[|X_i|]<\infty\) for all \(i\).
\pause
The \emph{mean vector} of \(X\) is defined as
\[
\mathbb{E}[X]
=
\begin{pmatrix}
\mathbb{E}[X_1]\\
\vdots\\
\mathbb{E}[X_d]
\end{pmatrix}
\in\mathbb{R}^d.
\vspace{-5mm}
\]
\end{definitionblock}
\vspace{-5mm}
\pause
\begin{itemize}
  \item The expectation is taken \emph{componentwise}.
  \item The mean vector represents the average location of the distribution.
  \item In finance, it corresponds to expected returns.
\end{itemize}
\pause
\medskip

\textbf{Linearity:}
for any matrix \(A\) and vector \(b\),
\[
\mathbb{E}[AX+b] = A\,\mathbb{E}[X] + b.
\]
\end{frame}

\begin{frame}{Example: Mean Vector}

\textbf{Example — Two-asset portfolio}

Let \(R=(R_1,R_2)\) be the vector of one-period returns of two assets.

Assume:
\[
\mathbb{E}[R_1]=5\%,
\qquad
\mathbb{E}[R_2]=8\%.
\]
\pause
\medskip

Then the mean vector is
\[
\mathbb{E}[R]
=
\begin{pmatrix}
0.05\\
0.08
\end{pmatrix}.
\]
\pause
\medskip

\textbf{Expected portfolio return.}
For portfolio weights \(w=(w_1,w_2)\),
\[
\mathbb{E}[w^\top R]
=
w^\top \mathbb{E}[R]
=
0.05\,w_1 + 0.08\,w_2.
\]
\pause
\medskip

\textbf{Interpretation:}
the mean vector summarizes the expected performance of each asset.

\end{frame}


\begin{frame}{Covariance}
\begin{definitionblock}{Covariance}
For \(X,Y\in L^2\),
\[
\mathrm{Cov}(X,Y)
= \mathbb{E}[(X-\mathbb{E}[X])(Y-\mathbb{E}[Y])]
= \mathbb{E}[XY]-\mathbb{E}[X]\mathbb{E}[Y].
\]\pause
\end{definitionblock}

\begin{itemize}
  \item \(\mathrm{Var}(X)=\mathrm{Cov}(X,X)\).\pause
  \item If \(X\) and \(Y\) are independent, then \(\mathrm{Cov}(X,Y)=0\).\pause
  \item
  \[
    \mathrm{Var}(X+Y)=\mathrm{Var}(X)+\mathrm{Var}(Y)+2\mathrm{Cov}(X,Y).
  \]\pause
\end{itemize}
\end{frame}

\begin{frame}{Covariance Matrix}
\begin{definitionblock}{Covariance Matrix}
For \(X=(X_1,\dots,X_d)\),
\[
\Sigma = \mathrm{Cov}(X)
= \big(\mathrm{Cov}(X_i,X_j)\big)_{i,j}.
\]\pause
\end{definitionblock}

\begin{itemize}
  \item \(\Sigma\) is symmetric and positive semi-definite.\pause
  \item Compact form:
  \[
    \Sigma = \mathbb{E}\big[(X-\mathbb{E}[X])(X-\mathbb{E}[X])^\top\big].
  \]\pause
\end{itemize}
\end{frame}

\begin{frame}{Linear Transformations}
Let \(X\in\mathbb{R}^d\), \(M\in\mathbb{R}^{p\times d}\), \(a\in\mathbb{R}^p\).
\pause
\begin{itemize}
  \item Mean:
  \[
    \mathbb{E}[MX+a]=M\mathbb{E}[X]+a.
  \]\pause
  \item Covariance:
  \[
    \mathrm{Cov}(MX+a)=M\,\mathrm{Cov}(X)\,M^\top.
  \]
\end{itemize}

\end{frame}

\begin{frame}{Portfolio Variance}

Let \(R=(R_1,\dots,R_d)\) be the vector of asset returns with
mean \(\mu=\mathbb{E}[R]\) and covariance matrix \(\Sigma=\mathrm{Cov}(R)\).

\medskip

For a portfolio with weights \(w=(w_1,\dots,w_d)\),
the portfolio return is
\[
R_p = w^\top R.
\]
\pause
\medskip

\begin{itemize}
  \item \textbf{Expected return:}
  \[
    \mathbb{E}[R_p] = w^\top \mu.
  \]
\pause
  \item \textbf{Variance:}
  \[
    \mathrm{Var}(R_p) = w^\top \Sigma\, w.
  \]\pause
\end{itemize}

\medskip

\textbf{Interpretation:}
\begin{itemize}
  \item Diagonal terms: individual asset risks.
  \item Off-diagonal terms: dependence and diversification effects.
\end{itemize}

\end{frame}

\begin{frame}{Correlation}
\begin{definitionblock}{Correlation}
For \(X,Y\in L^2\) with positive variances,
\[
\rho_{X,Y}
= \frac{\mathrm{Cov}(X,Y)}{\sqrt{\mathrm{Var}(X)\mathrm{Var}(Y)}} \in [-1,1].
\]
\end{definitionblock}
\pause
\begin{itemize}
  \item \(\rho=0\): no linear dependence.
  \item \(|\rho|=1\): perfect linear relationship.
  \item Correlation is invariant under affine transformations.
\end{itemize}
\end{frame}

% ======================================
% SUBSECTION — Gaussian Random Vectors
% ======================================
\subsection{Gaussian Random Vectors}

\begin{frame}{Gaussian Random Vectors}
\begin{definitionblock}{Gaussian Random Vector}
A random vector \(X\in\mathbb{R}^d\) is called \emph{Gaussian} (or \emph{multivariate normal})
if any linear combination
\[
a^\top X = \sum_{i=1}^d a_i X_i,
\quad a\in\mathbb{R}^d,
\]
is a (univariate) normal random variable.\pause
\end{definitionblock}

\begin{itemize}
  \item Denoted \(X\sim\mathcal{N}(\mu,\Sigma)\).\pause
  \item \(\mu=\mathbb{E}[X]\in\mathbb{R}^d\): mean vector.\pause
  \item \(\Sigma=\mathrm{Cov}(X)\in\mathbb{R}^{d\times d}\): covariance matrix.\pause
\end{itemize}

\end{frame}

\begin{frame}{Example and Counterexample}

\textbf{Example — Gaussian random vector}

Let \(Z=(Z_1,Z_2)\) where \(Z_1,Z_2\) are independent and
$
Z_1 \sim \mathcal{N}(0,1),
\,
Z_2 \sim \mathcal{N}(0,1).
$
\pause
Then
\[
Z \sim \mathcal{N}\!\left(
\begin{pmatrix}0\\0\end{pmatrix},
\begin{pmatrix}1&0\\0&1\end{pmatrix}
\right).
\]
\pause
For \(a\in\mathbb{R}^2\), \(a^\top Z\) is a combination of independent
normal variables, hence normal.
\vspace{-5mm}
\pause
\medskip

\textbf{Counterexample — Normal marginals but not Gaussian}

Let \(X\sim \mathcal{N}(0,1)\) and define
\[
Y =
\begin{cases}
X, & \text{with probability } \tfrac12,\\
-\,X, & \text{with probability } \tfrac12,
\end{cases}
\,\,\text{independently of \(X\).}
\vspace{-5mm}
\]
\pause
\begin{itemize}
  \item \(X\sim \mathcal{N}(0,1)\) and \(Y\sim \mathcal{N}(0,1)\),
  \item but the vector \((X,Y)\) is \emph{not} Gaussian,
  \item some linear combinations of \((X,Y)\) are not normally distributed.
\end{itemize}
\end{frame}

% --------------------------------------

\begin{frame}{Joint Density of a Gaussian Vector}

If \(\Sigma\) is positive definite, \(X\sim\mathcal{N}(\mu,\Sigma)\) admits the density
\[
f_X(x)
=
\frac{1}{(2\pi)^{d/2}|\Sigma|^{1/2}}
\exp\!\left(
-\frac12 (x-\mu)^\top \Sigma^{-1} (x-\mu)
\right),
\quad x\in\mathbb{R}^d.
\]
\pause
\begin{itemize}
  \item Elliptical level sets centered at \(\mu\).
  \item Shape determined by \(\Sigma\).
\end{itemize}

\end{frame}

% --------------------------------------

\begin{frame}{Key Properties of Gaussian Vectors}

Let \(X\sim\mathcal{N}(\mu,\Sigma)\).
\pause
\begin{itemize}
  \item Any subvector of \(X\) is Gaussian.\pause
  \item Any linear transformation is Gaussian:
  \[
    AX+b \sim \mathcal{N}(A\mu+b,\;A\Sigma A^\top).
  \]\pause
  \item A Gaussian vector is fully characterized by \((\mu,\Sigma)\).\pause
\end{itemize}

\medskip

\textbf{Important consequence:}
for Gaussian vectors,
\[
\text{zero covariance} \;\Longleftrightarrow\; \text{independence}.
\]

\end{frame}

% --------------------------------------

\begin{frame}{Marginals and Independence}

Let \(X=(X_1,\dots,X_d)\sim\mathcal{N}(\mu,\Sigma)\).

\begin{itemize}
  \item Each component \(X_i\sim\mathcal{N}(\mu_i,\Sigma_{ii})\).\pause
  \item Any subvector is multivariate normal.\pause
  \item If \(\Sigma_{ij}=0\), then \(X_i\) and \(X_j\) are independent.\pause
\end{itemize}

\medskip

\textbf{Contrast:}
this equivalence between independence and zero covariance
is \emph{specific} to Gaussian vectors.

\end{frame}

% --------------------------------------

\begin{frame}{Gaussian Vectors in Finance}

\begin{itemize}
  \item Asset returns often modeled as:
  \[
    R \sim \mathcal{N}(\mu,\Sigma).
  \]\pause
  \item Portfolio return:
  \[
    w^\top R \sim \mathcal{N}(w^\top\mu,\; w^\top\Sigma w).
  \]\pause
  \item Mean–variance analysis relies entirely on \((\mu,\Sigma)\).
\end{itemize}

\medskip

\textbf{Limitations:}
\begin{itemize}
  \item No skewness.
  \item Thin tails (underestimates extreme risks).
\end{itemize}

\end{frame}

% --------------------------------------

\begin{frame}{Summary: Gaussian Random Vectors}

\begin{itemize}
  \item Gaussian vectors generalize the normal distribution to \(\mathbb{R}^d\).\pause
  \item Fully characterized by mean vector and covariance matrix.\pause
  \item Linear transformations and conditional laws remain Gaussian.\pause
  \item Central tool in portfolio theory and multivariate statistics.
\end{itemize}

\end{frame}
% ======================================
% SECTION — Convergence of Random Variables
% ======================================
\section{Convergence of Random Variables}

\begin{frame}{Outline}
  \tableofcontents[currentsection, hideothersubsections]
\end{frame}

% --------------------------------------
\subsection{Why Convergence?}

\begin{frame}{Why Do We Need Convergence?}

In probability and statistics, we often work with sequences of random variables
\[
(X_n)_{n\ge1}.
\]
\pause
\medskip

\begin{itemize}
  \item Approximating a complicated quantity by a simpler one.\pause
  \item Studying large-sample behavior (statistics, econometrics).\pause
  \item Understanding long-run averages and fluctuations.\pause
\end{itemize}

\medskip

\textbf{Key question:}
In what sense does \(X_n\) get close to a limiting random variable \(X\)?
\pause
\end{frame}

% --------------------------------------
\subsection{Different Types of Convergence}

\begin{frame}{Different Notions of Convergence}

Several notions of convergence coexist in probability theory:
\pause
\begin{itemize}
  \item almost sure convergence,\pause
  \item convergence in probability,\pause
  \item convergence in \(L^p\),\pause
  \item convergence in distribution (in law).\pause
\end{itemize}

\medskip

They differ in strength and interpretation, but are closely related.

\end{frame}

% --------------------------------------
\begin{frame}{Convergence in Probability}

\begin{definitionblock}{Convergence in Probability}
A sequence \((X_n)_{n\ge1}\) converges to \(X\) in probability,
denoted
\[
X_n \xrightarrow{\mathbb{P}} X,
\]
if for every \(\varepsilon>0\),
\[
\lim_{n\to\infty}\mathbb{P}(|X_n-X|>\varepsilon)=0.
\]
\end{definitionblock}
\pause
\begin{itemize}
  \item Deviations larger than \(\varepsilon\) become unlikely.
  \item Central notion in statistics and estimation.
\end{itemize}

\end{frame}

\begin{frame}{Example: Convergence in Probability}

\textbf{Example (Sample mean).}

Let \(X_1,X_2,\dots\) be i.i.d.\ random variables with
\[
\mathbb{E}[X_1]=\mu,
\qquad
\mathrm{Var}(X_1)=\sigma^2<\infty.
\]
\pause
Define the sample mean
$
\bar X_n=\frac1n\sum_{k=1}^n X_k.
$
\pause
\medskip
Then, for any \(\varepsilon>0\),
\[
\mathbb{P}\big(|\bar X_n-\mu|>\varepsilon\big)
\;\le\;
\frac{\mathrm{Var}(\bar X_n)}{\varepsilon^2}
=
\frac{\sigma^2}{n\varepsilon^2}
\;\xrightarrow[n\to\infty]{}\;0.
\]
\pause
\medskip

\textbf{Conclusion:}
\[
\bar X_n \xrightarrow{\mathbb{P}} \mu.
\]

\medskip

This is the \emph{Weak Law of Large Numbers}.

\end{frame}

% --------------------------------------
\begin{frame}{Almost Sure Convergence}

\begin{definitionblock}{Almost Sure Convergence}
A sequence \((X_n)_{n\ge1}\) converges to \(X\) almost surely,
denoted
\[
X_n \xrightarrow{\text{a.s.}} X,
\]
if
\[
\mathbb{P}\big(\omega : X_n(\omega)\to X(\omega)\big)=1.
\]\pause
\end{definitionblock}

\begin{itemize}
  \item Pointwise convergence for almost all outcomes.
  \item Strongest and most intuitive notion.
\end{itemize}

\end{frame}

\begin{frame}{Example: Almost Sure Convergence}

\textbf{Example (Strong Law of Large Numbers).}

Let \(X_1,X_2,\dots\) be i.i.d.\ random variables such that
\[
\mathbb{E}[|X_1|] < \infty.
\]
\pause
Define the sample mean
\[
\bar X_n=\frac1n\sum_{k=1}^n X_k.
\]

\medskip

Then
\[
\mathbb{P}\Big(
\omega:\ \lim_{n\to\infty} \bar X_n(\omega)=\mathbb{E}[X_1]
\Big)=1.
\]
\pause
\medskip

\textbf{Conclusion:}
\[
\bar X_n \xrightarrow{\text{a.s.}} \mathbb{E}[X_1].
\]

\medskip

\textbf{Interpretation:}
for almost every outcome \(\omega\), the empirical average
stabilizes and converges to the true mean.

\end{frame}

% --------------------------------------
\begin{frame}{Convergence in \texorpdfstring{$L^p$}{Lp}}

\begin{definitionblock}{\(L^p\) Convergence}
For \(p\ge1\), we say that \(X_n\) converges to \(X\) in \(L^p\),
denoted
\[
X_n \xrightarrow{L^p} X,
\]
if
\[
\lim_{n\to\infty}\mathbb{E}[|X_n-X|^p]=0.
\]\pause
\end{definitionblock}

\begin{itemize}
  \item Controls the mean size of the error.
  \item \(L^2\) convergence is especially important (variance).
\end{itemize}

\end{frame}

\begin{frame}{Example: Convergence in \texorpdfstring{$L^p$}{Lp}}

\textbf{Example (Sample mean in \(L^2\)).}

Let \(X_1,X_2,\dots\) be i.i.d.\ random variables with
\[
\mathbb{E}[X_1]=\mu,
\qquad
\mathrm{Var}(X_1)=\sigma^2<\infty.
\]
\pause
Define
\[
\bar X_n=\frac1n\sum_{k=1}^n X_k.
\]

\medskip

Then
\[
\mathbb{E}\big[(\bar X_n-\mu)^2\big]
=
\mathrm{Var}(\bar X_n)
=
\frac{\sigma^2}{n}
\;\xrightarrow[n\to\infty]{}\;0.
\]

\medskip

\textbf{Conclusion:}
$\bar X_n \xrightarrow{L^2} \mu.$

\medskip

\textbf{Interpretation:}
the mean squared error between \(\bar X_n\) and \(\mu\) goes to zero.

\end{frame}

% --------------------------------------
\begin{frame}{Convergence in Distribution}

\begin{definitionblock}{Convergence in Distribution}
We say that \(X_n\) converges in distribution to \(X\),
denoted
\[
X_n \xrightarrow{d} X,
\]
if for every bounded continuous function \(\varphi\),
\[
\lim_{n\to\infty}\mathbb{E}[\varphi(X_n)]
=
\mathbb{E}[\varphi(X)].
\]\pause
\end{definitionblock}

\begin{itemize}
  \item Concerns the convergence of \emph{laws}.
  \item Random variables need not live on the same probability space.
\end{itemize}

\end{frame}

\begin{frame}{Example: Convergence in Distribution}

\textbf{Example (Central Limit Theorem).}

Let \(X_1,X_2,\dots\) be i.i.d.\ random variables with
\[
\mathbb{E}[X_1]=\mu,
\qquad
\mathrm{Var}(X_1)=\sigma^2<\infty.
\]

Define the normalized sum
\[
Z_n
=
\frac{\sqrt{n}\,(\bar X_n-\mu)}{\sigma}
=
\frac{1}{\sigma\sqrt{n}}
\sum_{k=1}^n (X_k-\mu).
\]

\medskip

Then
$
Z_n \xrightarrow{d} \mathcal{N}(0,1).
$

\medskip

\textbf{Interpretation:}
although \(Z_n\) does not converge to a fixed random variable,
its \emph{distribution} converges to the standard normal law.

\medskip

This is the cornerstone of asymptotic statistics.

\end{frame}

% \begin{frame}{CLT: A Visual Picture (Standardized Sums)}
% \textbf{Concrete example.} Let \(X_k\in\{-1,+1\}\) with \(\mathbb{P}(X_k=\pm1)=\tfrac12\).
% Then \(\mu=0\), \(\sigma^2=1\), and
% \[
% Z_n=\frac{1}{\sqrt{n}}\sum_{k=1}^n X_k.
% \]

% \begin{center}
% \begin{tikzpicture}
% \begin{axis}[
%   width=11.2cm, height=6.2cm,
%   xmin=-3.2, xmax=3.2,
%   ymin=0, ymax=0.55,
%   axis lines=left,
%   xlabel={$z$}, ylabel={density / pmf},
%   legend style={at={(0.02,0.98)},anchor=north west, draw=none, fill=none},
% ]

% % Standard normal density
% \addplot[thick, domain=-3.2:3.2, samples=200]
%   {1/sqrt(2*pi)*exp(-x^2/2)};
% \addlegendentry{$\mathcal{N}(0,1)$}

% % Z_1
% \addplot+[ycomb, mark=*, mark size=1.0pt]
%   coordinates {(-1,0.5) (1,0.5)};
% \addlegendentry{$Z_1$}

% % Z_5
% \addplot+[ycomb, mark=*, mark size=0.7pt, opacity=0.75]
% coordinates {
%   (-2.2360679,0.03125)
%   (-1.3416408,0.15625)
%   (-0.4472136,0.3125)
%   ( 0.4472136,0.3125)
%   ( 1.3416408,0.15625)
%   ( 2.2360679,0.03125)
% };
% \addlegendentry{$Z_5$}

% \end{axis}
% \end{tikzpicture}
% \end{center}

% \vspace{-2mm}
% \begin{itemize}
%   \item The law is discrete for finite \(n\) (spikes).
%   \item As \(n\) grows, it aligns with the smooth \(\mathcal{N}(0,1)\) curve: \(Z_n \xrightarrow{d}\mathcal{N}(0,1)\).
% \end{itemize}
% \end{frame}

% --------------------------------------
\subsection{Relations Between Convergences}

\begin{frame}{Relations Between the Notions}

The following implications hold in general:
\[
\text{almost surely}
\;\Rightarrow\;
\text{in probability}
\;\Rightarrow\;
\text{in distribution}.
\]
\pause
Moreover, for any \(p\ge1\),
\[
L^p
\;\Rightarrow\;
\text{in probability}.
\vspace{-5mm}
\]
\pause
\medskip

\begin{itemize}
  \item None of the converse implications hold in general.
  \item \(L^p\) convergence does \emph{not} imply almost sure convergence.
\end{itemize}
\pause
\begin{definitionblock}{Special Case}
If \(X_n\xrightarrow{d} c\) where \(c\) is a constant,
then
\[
X_n\xrightarrow{\mathbb{P}} c.
\]
\end{definitionblock}

\end{frame}

\begin{frame}{Example: Almost Sure $\Rightarrow$ Probability $\Rightarrow$ Distribution}

Let \(X_n(\omega)=\frac{1}{n}\) for all \(\omega\in\Omega\).
\pause
\medskip

Then:
\begin{itemize}
  \item For every \(\omega\),
  \[
    X_n(\omega)\to 0,
  \]
  hence
  \[
    X_n \xrightarrow{\text{a.s.}} 0.
  \]
  \item Therefore,
  \[
    X_n \xrightarrow{\mathbb{P}} 0
    \quad\text{and}\quad
    X_n \xrightarrow{d} 0.
  \]
\end{itemize}

\medskip

\textbf{Conclusion:} almost sure convergence implies all weaker notions.

\end{frame}

\begin{frame}{Counterexample: In Probability but not Almost Surely}

Let \((U_n)_{n\ge1}\) be i.i.d.\ uniform on \([0,1]\), and define
\[
X_n = \mathbf{1}_{\{U_n \le 1/n\}}.
\]
\pause
Then:
\begin{itemize}
  \item \(\mathbb{P}(X_n=1)=\frac{1}{n}\to0\), hence
  \[
    X_n \xrightarrow{\mathbb{P}} 0.
  \]
  \item However,
  \[
    \sum_{n=1}^\infty \mathbb{P}(X_n=1)=\infty,
  \]
  so by Borel--Cantelli,
  \[
    \mathbb{P}(X_n=1 \text{ infinitely often})=1.
  \]
\end{itemize}
\pause
\medskip

\textbf{Conclusion:}
\[
X_n \not\xrightarrow{\text{a.s.}} 0.
\]

\end{frame}

\begin{frame}{Counterexample: In Distribution but not in Probability}

Let \(X_n\sim \mathcal{N}(0,1)\) for all \(n\), and assume the \(X_n\) are independent.
\pause
\medskip

Then:
\begin{itemize}
  \item The distribution does not depend on \(n\), so
  \[
    X_n \xrightarrow{d} \mathcal{N}(0,1).
  \]
  \item But for any random variable \(X\),
  \[
    \mathbb{P}(|X_n-X|>\varepsilon) \not\to 0
  \]
  in general.
\end{itemize}
\pause
\medskip

\textbf{Conclusion:}
convergence in distribution does not imply convergence in probability
unless the limit is a constant.

\end{frame}

% --------------------------------------
\subsection{Preservation of Convergence}

\begin{frame}{Stability Properties}

\begin{itemize}
  \item If \(X_n\xrightarrow{\mathbb{P}} X\) and
        \(Y_n\xrightarrow{\mathbb{P}} Y\), then
  \[
  X_n+Y_n \xrightarrow{\mathbb{P}} X+Y,
  \qquad
  X_nY_n \xrightarrow{\mathbb{P}} XY.
  \]\pause
  \item Same results hold for almost sure convergence.
\end{itemize}\pause

\medskip

\begin{definitionblock}{Continuous Mapping Theorem}
If \(X_n\xrightarrow{\mathbb{P}} X\) and \(g\) is continuous,
then
\[
g(X_n)\xrightarrow{\mathbb{P}} g(X).
\]
\end{definitionblock}

\end{frame}

% --------------------------------------
\begin{frame}{Slutsky's Lemma}

\begin{definitionblock}{Slutsky's Lemma}
If
\[
X_n\xrightarrow{d} X,
\qquad
Y_n\xrightarrow{\mathbb{P}} c,
\]
then
\[
(X_n,Y_n)\xrightarrow{d}(X,c).
\]\pause
\end{definitionblock}

\medskip

In particular:
\[
X_n+Y_n\xrightarrow{d} X+c,
\quad
X_nY_n\xrightarrow{d} Xc.
\]

\end{frame}

\begin{frame}{Application: Continuous Mapping and Stability}

\textbf{Example (Transformation of an estimator).}

Let \(X_1,X_2,\dots\) be i.i.d.\ with mean \(\mu\), and define
\[
\bar X_n = \frac1n\sum_{k=1}^n X_k.
\]

By the Law of Large Numbers,
\[
\bar X_n \xrightarrow{\mathbb{P}} \mu.
\]
\pause
\medskip

Let \(g(x)=x^2\), which is continuous.  
By the Continuous Mapping Theorem,
\[
(\bar X_n)^2 \xrightarrow{\mathbb{P}} \mu^2.
\]

\medskip
\pause
\textbf{Interpretation:}
plug-in estimators preserve convergence in probability
under continuous transformations.

\end{frame}

\begin{frame}{Application: Slutsky's Lemma}

\textbf{Example (Studentized mean).}

Let \(X_1,X_2,\dots\) be i.i.d.\ with
$
\mathbb{E}[X_1]=\mu,
\qquad
\mathrm{Var}(X_1)=\sigma^2.
$
Define:
\[
\bar X_n=\frac1n\sum_{k=1}^n X_k,
\qquad
S_n^2=\frac1n\sum_{k=1}^n (X_k-\bar X_n)^2.
\]

\medskip

Then:
\[
\sqrt{n}(\bar X_n-\mu) \xrightarrow{d} \mathcal{N}(0,\sigma^2),
\qquad
S_n \xrightarrow{\mathbb{P}} \sigma.
\]

\medskip

By Slutsky's lemma,
\[
\frac{\sqrt{n}(\bar X_n-\mu)}{S_n}
\xrightarrow{d}
\mathcal{N}(0,1).
\]

\medskip

\textbf{Interpretation:}
unknown parameters can be replaced by consistent estimators.

\end{frame}

% --------------------------------------
\subsection{Law of Large Numbers}

\begin{frame}{Law of Large Numbers}

Let \(X_1,X_2,\dots\) be i.i.d.\ random variables with
\(\mathbb{E}[X_1]=\mu\).
\pause
\medskip

\begin{definitionblock}{Weak Law of Large Numbers}
If \(\mathrm{Var}(X_1)<\infty\), then
\[
\bar X_n = \frac1n\sum_{k=1}^n X_k
\xrightarrow{L^2} \mu.
\]\pause
\end{definitionblock}

\begin{definitionblock}{Strong Law of Large Numbers}
If \(\mathbb{E}[|X_1|]<\infty\), then
\[
\bar X_n \xrightarrow{\text{a.s.}} \mu.
\]\pause
\end{definitionblock}

\end{frame}

% --------------------------------------
\subsection{Central Limit Theorem}

\begin{frame}{Central Limit Theorem}

Let \(X_1,X_2,\dots\) be i.i.d.\ with mean \(\mu\) and variance \(\sigma^2>0\).
\pause
\begin{definitionblock}{Central Limit Theorem}
\[
\sqrt{n}\,(\bar X_n-\mu)
\xrightarrow{d}
\mathcal{N}(0,\sigma^2).
\]
\end{definitionblock}

\medskip

Equivalently,
\[
\frac{\bar X_n-\mu}{\sigma/\sqrt{n}}
\xrightarrow{d}
\mathcal{N}(0,1).
\]

\end{frame}

% --------------------------------------
\begin{frame}{Delta Method}

\begin{definitionblock}{Delta Method}
Let \(g:\mathbb{R}\to\mathbb{R}\) be continuously differentiable.
Then
\[
\sqrt{n}\big(g(\bar X_n)-g(\mu)\big)
\xrightarrow{d}
\mathcal{N}\big(0,(g'(\mu))^2\sigma^2\big).
\]
\end{definitionblock}
\pause
\begin{itemize}
  \item Allows nonlinear transformations.
  \item Widely used in statistics and econometrics.
\end{itemize}

\end{frame}

% --------------------------------------


\section{References}
\begin{frame}{Suggested References}
\begin{itemize}
  \item Carl P. Simon and Lawrence Blume, \textit{Mathematics for Economists}.
  \item J. Douglas Carroll and Paul E. Green, \textit{Mathematical Tools for Applied Multivariate Analysis}.
  \item (Optional, more advanced) Bruce Hajek, \textit{Random Processes}; David Williams, \textit{Probability with Martingales}.
\end{itemize}
\end{frame}

\end{document}